\begin{figure}[htbp]
\centering
\begin{tikzpicture}[node distance=2mm,timelinebox/.style={rounded corners=2mm,draw=DeepBlue,thick,text width=.24\textwidth,minimum height=.9in,align=center},arrow/.style={-{Latex[length=2mm]},thick,draw=PrimaryRed}]
\node[timelinebox,fill=SoftBlue] (a) {\textbf{0--2 min}\\Preview first/next/last\\Introduce autism respectfully};
\node[timelinebox,fill=SoftGreen,right=of a] (b) {\textbf{2--4 min}\\Normalize communication\\and sensory differences};
\node[timelinebox,fill=SoftGold,right=of b] (c) {\textbf{4--8 min}\\Teach five peer actions\\Practice a coding scenario};
\node[timelinebox,fill=SoftLavender,right=of c] (d) {\textbf{8--10 min}\\Choose one action\\Explain why it supports belonging};
\draw[arrow] (a)--(b); \draw[arrow] (b)--(c); \draw[arrow] (c)--(d);
\end{tikzpicture}
\caption{The 10-minute lesson sequence. The visual timeline reduces uncertainty and connects each teaching move to its purpose.}
\label{fig:lesson-timeline}
\end{figure}